
\documentclass[a4paper, 11pt]{article}

    %%%%%%%%%%%%%%%%%%%%%%%%%%%%%%
    \oddsidemargin  0.0in
    \evensidemargin 0.0in
    \textwidth      6.2in
    \headheight     0.0in
    \topmargin      0.0in
    \textheight=9.0in
    \parindent=0in
    %%%%%%%%%%%%%%%%%%%%%%%%%%%%%%


\usepackage{amsthm,amsfonts,amssymb,amsmath}
\usepackage{graphicx}   %To include graphics
\usepackage[small,compact]{titlesec}
\usepackage{verbatim} %To write multi-line comments easily


\usepackage[utf8]{inputenc} 
\usepackage[francais]{babel}


\begin{document}

\begin{center}
{\bf STT-2100 Régression }\\
\end{center}
\begin{center}
{\bf Examen 1 - 18 octobre 2022 }\\
\end{center}


\vspace{0.2 in}
Cet examen dure 110 minutes. \\

Vous avez droit à une calculatrice parmi la liste des calculatrices acceptées par la FSG. Si les valeurs exactes que vous recherchez ne sont pas dans la table, prendre la valeur la plus proche.\\

Vous pouvez utiliser le verso de l'examen comme papier brouillon.

\vspace{0.4 in}

{\bf Matricule:} \quad\rule{2in}{.01in}

\vspace{.5in}

{\bf Nom:}\quad\rule{2in}{.01in}{\bf Prénom:}\quad\rule{2in}{.01in}\\

%{\bf Your Section Name:}\quad\rule{2in}{.01in}\\

{\bf Signature:}\quad\rule{3in}{.01in}\\

\sf {


\vspace{0.3in}

\begin{center}
{\def\arraystretch{1.5}\tiny\normalsize
\begin{tabular}[t]{|l|r|r|}
\hline
Question & Total & Correct \\
\hline
\hline
PARTIE 1     &  20   &      \\\hline
PARTIE 2     &  25    &       \\ \hline             
PARTIE 3 - Q1   &  20   &       \\ \hline 
PARTIE 3 - Q2     &  12   &       \\ \hline
PARTIE 3 - Q3   &  23   &       \\   \hline                   
%7     &      &       \\ \hline 
Total   & 100  &                \\ \hline
\end{tabular}
}

\vspace{1 in}
{\bf{ *** Attendre le signal avant de tourner cette page.***}}

\vspace{0.5 in }

{\bf ASSUREZ-VOUS D'AVOIR LES PAGES 1 à 9}

\vspace{0.5 in}

\end{center}

%%%%%%%%%%%%%%%%%%%%%%%%%%%%%%%%%%%%%%%%%%%%%%%%%%%%%%%%%%%%%%%%%%%%%%%%%%%%%%
\newpage
%{\bf Your ANDREW ID:} \IDboxes
%\vspace{.2in}


\section*{PARTIE 1 - VRAI OU FAUX?}

\textbf{Pour chaque énoncé ci-dessous indiquez s'il est vrai ou faux. \underline{Justifiez} brièvement votre réponse. \\
Chaque question vaut 4 points. Aucun point ne sera accordé pour une réponse sans justification. }
\begin{enumerate}
\item  Si la valeur-p au test F global pour la régression linéaire est 0,30, alors je sais que l'intervalle de
confiance à 95\% pour $\beta_1$ ne contient pas la valeur 0.
\vspace{1.2 in}
\item Si je fais l'hypothèse d’un modèle linéaire dans lequel les erreurs ont une distribution uniforme entre -1 et 1 au lieu d’une distribution normale, les estimateurs des moindres carrés sont quand même sans biais pour les coefficients de la régression.
\vspace{1.2 in}
\item La covariance entre $\beta_0$ et $\beta_1$ est nulle. 
\vspace{1.2 in}
\item L'intervalle de confiance à $100(1-\alpha)\%$ pour la valeur moyenne $Y_0 = \beta_0 + \beta_1 x_0$ est toujours plus court que l'intervalle de prévision à $100(1-\alpha)\%$ pour la valeur d'une observation au point $x_0$. 
\vspace{1.2 in}
\item La somme des carrés de l’erreur résiduelle (SCres) est égale à l'estimateur de la variance des termes d’erreur ($\widehat{\sigma}^2$).

\vspace{1.2 in}
%



%


%
\end{enumerate}


\newpage

% SHORT-ANSWER QUESTIONS

\vspace*{-0.7 in}
\section*{PARTIE 2 - RÉPONSES COURTES (5 pts par question)}
\vspace{0.1 in}
\begin{enumerate}
\item Pourquoi voudrait-on utiliser les résidus standardisés au lieu des résidus bruts pour tester la validité des postulats?
\vspace{1.2 in }
\item Tracez un graphique des résidus standardisés en fonction des valeurs prédites pour des observations qui ne satisfont pas l'hypothèse d'homoscédasticité. Supposez que $\hat{\beta}_1$ est négatif.
\vspace{1.5 in }
\item Pourquoi l'hypothèse de normalité des erreurs est-elle importante dans la régression linéaire simple ?
\vspace{1.2 in }
\item Dans la planification de la récolte des données, nous pouvons parfois choisir pour quelles valeurs de $x$ nous prenons des observations. Quel est l'impact du choix des valeurs de $x$ sur la précision des estimateurs du modèle de régression linéaire ?
\vspace{1.2 in }
\item Supposons que la vraie relation entre deux variables est $Y_i = x_i^2 + \epsilon_i$ avec $\epsilon_i \overset{iid}{\sim} N(0,\sigma^2)$, mais qu'on ajuste un modèle de régression linéaire simple entre $x$ et $Y$. Tracez le graphique des résidus standardisés en fonction des valeurs prédites que vous observeriez dans ce cas. 





\end{enumerate}


\newpage 

%LONG QUESTIONS

\section*{PARTIE 3}

\begin{enumerate}


\item 
Soit un jeu de données de taille $n$ auquel on ajuste le modèle de régression linéaire suivant:
\begin{equation*}
Y_i = \beta_1 x_i + \varepsilon_i, \mbox{ pour } i = 1,\ldots,n
\end{equation*}
sous l'hypothèse que $\varepsilon_i \overset{iid}{\sim}N(0, \sigma^2)$. On peut montrer que l'estimateur des moindres carrés est $$\widehat{\beta}_1 = \frac{\displaystyle \sum_{i=1}^n(x_iy_i)}{\displaystyle \sum_{i=1}^n x_i^2}.$$   
\begin{enumerate}
    \item (3 points) Donnez la distribution, l'espérance et la variance de $Y_i$? 
    \vspace{1.4 in }
    \item (6 points) Calculez l'espérance de $\widehat{\beta}_1$ et donnez l'expression de biais pour $\beta_1$. Justifiez toutes les étapes de votre démarche.
    \newpage
    \item (8 points) Calculez la variance de $\widehat{\beta}_1$. Justifiez vos calculs et toutes les étapes de votre démarche.
    \vspace{5 in }
    \item (3 points) Donnez la distribution de $\widehat{\beta}_1$.
\end{enumerate}

\newpage

\item 
Dans le Guide de consommation de carburant 2006 publié par le Ministère des Ressources
naturelles du Canada, on retrouve des données sur la cylindrée et la consommation d'essence
en ville de plusieurs modèles de véhicules. Vous trouverez à la page suivante une partie de la sortie SAS obtenue en ajustant un mod\`ele de r\'egression lin\'eaire simple \`a ce jeu de données. Certaines des entr\'ees du tableau ont \'et\'e masqu\'ees par un ? et d'autres par des lettres majuscules (AAA,BBB,CCC,DDD). Trouvez les valeurs des entrées données par une lettre majuscule et \underline{justifiez} vos r\'eponses. \\

Chaque sous-question vaut 3 points.\\

\vspace{0.1 in}

\begin{enumerate}
\item[]AAA-
\vspace{1 in}
\item[]BBB-
\vspace{1 in}
\item[]CCC-
\vspace{1 in}
\item[]DDD-
\end{enumerate}

\newpage 
 \begin{verbatim}
                  Procédure REG
                 Modèle : MODEL1
                 Variable dépendante : consom
                 
 Nombre d'observations lues                             20
 Nombre d'observations utilisées                        20


                    Analyse de variance

                              Somme des       Moyenne   Valeur
Source                 DDL       carrés   quadratique     de F   Pr > F

Modèle                   1    146.96404            ?       BBB       ?
Erreur                 AAA     27.75346      1.54186
Total sommes corrigées   ?            ?


Racine MSE   1.24172       R carré     CCC


                  Valeurs estimées des paramètres

                  Valeur estimée    Erreur
Variable    DDL   des paramètres      type    Valeur test t    Pr > |t|

Intercept     1         5.76226    0.72256             DDD       <.0001
cylind        1         2.03523          ?            9.76       <.0001



 \end{verbatim}
 
\newpage

\item On a noté pour plusieurs gâteaux la température de cuisson (en Farenheit) et la palatabilité du gâteau. La palatabilité est une mesure de goût: plus ce score est élevé, meilleur est le gâteau. Le résultat de l'analyse de régression linéaire simple est donné ci-dessous: 

\begin{verbatim}

              Estimate   Std. Error   t value   Pr(>|t|)  
(Intercept)  -26.70877     16.46779    -1.622      0.131  
temp           0.09639      0.04704     2.049      0.063 

Residual standard error: 1.33 on 12 degrees of freedom
Multiple R-squared: 0.2592,	 Adjusted R-squared: 0.1975 

\end{verbatim}
On note aussi que la moyenne des températures observées dans l'échantillon est de $350\,^{\circ}\mathrm{F}$. Pour chacune des questions suivantes, \underline{justifiez} votre réponse.
\vspace{0.1 in}

\begin{enumerate}
\item (3 pts) Au niveau $\alpha = 0,1$, y a-t-il une relation linéaire significative entre la température de cuisson et la palatabilité? Justifiez votre réponse. 
\vspace{1.4 in}
\item (4 pts) Votre client conclut qu'en augmentant de $10\,^{\circ}\mathrm{F}$ la température de cuisson, il augmentera la palatabilité de son gâteau d'environ une unité. Qu'en pensez-vous?
\vspace{1.4 in}
\item (3 pts) Quelle est la palatabilité prédite pour un gâteau cuit à $400\,^{\circ}\mathrm{F}$? 
\vspace{1.4 in}
\item (3 pts) Donnez une estimation non biaisé de la variance des erreurs.
\vspace{1 in}
\item (3 pts) Interprétez la valeur du coefficient de détermination.
\vspace{1.4 in}
\item (3 pts) Votre client conclut que si la température de son four est de $5000\,^{\circ}\mathrm{F}$, la palatabilité sera très élevée et il obtiendra un excellent gâteau. Quel est votre avis sur sa conclusion?
\vspace{1.5 in}
\item (4 pts) Donnez la valeur de $S_{xx}$.
\end{enumerate}



\end{document}

\item  %Question 2

On obtient les résultats numériques suivants :
\begin{eqnarray*}
\sum x_i=1872 &  \sum Y_i=32,3 & \sum x_i^2=294320 \\
\sum Y_i^2=93,11 & \sum x_iY_i=5156,2 & \hat{\beta}_1 = 0.0478\\
\end{eqnarray*}
\begin{enumerate}
\item Testez la pertinence de la régression au niveau $\alpha=1\%$.
\item Établissez l'intervalle de confiance de la valeur prédite pour $x=150$ cm au niveau de confiance $95\%$.

\item Votre nouveau collègue, Renaud, se présente à un 5 à 7 avec sa femme, Véronique. Renaud est agé de 53 ans, mais sa femme vCelle-ci semble n'avoir que  
\end{enumerate}

\end{enumerate}

\end{document}

BONUS : Pourquoi ne peut-on pas tester la linéarité directement à partir des variables $Y_i$? Après tout, elles devraient suivre une loi normale...




Deux collègues ont ajusté un modèle de régression linéaire simple pour prédire le montant des dommages à une maison en cas d'incendie à l'aide de la distance entre la maison et la station d'incendie la plus proche. Dans leur jeu de données, les distances entre la maison et la station d'incendie varient entre 200m et 5 km. Ils se questionnent sur la possibilité de prédire le montant des dommages en cas d'indencie pour une maison située à 10km de la station d'incendie la plus proche. 
Collègue 1 : Il est dangereux d'extrapoler ainsi avec notre modèle. 

Vous surprenez la discussion suivante entre deux collègues: 
Collègue 1 : J'ai ajusté un modèle de régression linéaire simple pour 
Collègue 2 : Pourrais-tu prédire le montant des dommages à ma maison en cas d'incendie. Elle est située à 8km d'une station d'incendie. 
Collègue 1 : Désolé, mais mon jeu de données avait des observations avec une distance entre 200 m et 5 km à une station d'incendie, donc je ne pense pas que ce soit approprié. 
Ce n'est pas grave, l'intervalle sera simplement plus long. 
